\section{DiBud: Budgeted Exact-Prefix Fusion}
\label{sec:method}
Given a query, a corpus snapshot, and a read budget $B$, DiBud reads the dense and sparse rankings on demand and returns a certified prefix of their complete RRF ranking. The budget determines how much can be read; certification determines how much can be returned.

\subsection{Budget and output}
Let $r_i(x)$ denote the one-based rank of document $x$ in channel $i\in\{D,S\}$, and let $g(r)=1/(c+r)$ for rank constant $c\geq 0$. Set $r_i(x)=\infty$ and $g(\infty)=0$ if $x$ is outside that channel's support. The complete fusion score is
\begin{equation}
 F(x)=\sum_{i\in\{D,S\}}g(r_i(x)).
 \label{eq:score}
\end{equation}
Sorting the union of the channel supports by decreasing $F(x)$, with a fixed document-identifier order to break ties, defines the complete sequence $T^\star$.

Both channels supply resumable exact ranking prefixes from the same snapshot. If $d_D$ and $d_S$ are their read depths and $A$ is the output sequence, the required contract is
\begin{equation}
 d_D+d_S\leq B,\qquad A=T^\star_{1:K},\quad K=|A|.
 \label{eq:contract}
\end{equation}
$K$ is not specified in advance. The output is empty if no position can be certified. The budget counts channel entries read, including separate reads of the same document in different channels.

\subsection{Prefix certification}
Let $S_i$ be the documents already read from channel $i$. An unread contribution in that channel is bounded by
\begin{equation}
 u_i=\begin{cases}
 0,&\text{if confirmed exhausted},\\
 g(d_i+1),&\text{otherwise}.
 \end{cases}
 \label{eq:unread}
\end{equation}
For an observed document $x$, the score bounds are
\begin{align}
 \ell(x)&=\sum_{i:x\in S_i}g(r_i(x)),\label{eq:lower}\\
 h(x)&=\ell(x)+\sum_{i:x\notin S_i}u_i.\label{eq:upper}
\end{align}
A document unseen in both channels has score at most $h_{\varnothing}=u_D+u_S$.

Among observed documents not yet emitted, select $x$ with the highest lower bound, breaking ties by the fixed identifier order. Certify $x$ as the next output when its lower bound exceeds every remaining competitor's upper bound and $h_{\varnothing}$. Equality with an observed competitor is allowed only when $x$ wins the identifier tie; comparison with the unseen bound remains strict. Append $x$ and repeat until the next position cannot be certified.

The bounds ensure $\ell(x)\leq F(x)\leq h(x)$. Each emitted document therefore precedes every remaining document, including any not yet observed. Applying this argument at each output position establishes the prefix equality in~\eqref{eq:contract}.

\subsection{Reading and stopping}
When the next position cannot be certified, let $y$ be the observed competitor with the highest upper bound, excluding $x$ and the emitted prefix and breaking ties by identifier. If $h(y)\geq h_{\varnothing}$ and $y$ is missing a contribution from exactly one channel, advance that channel. This either reveals the contribution or tightens its upper bound.

Otherwise, advance the channel whose next batch offers the larger reduction in the unread bound. For a batch of $b$ entries in a channel that remains open after the read, this reduction is
\begin{equation}
 \Delta_i=g(d_i+1)-g(d_i+b+1).
 \label{eq:drop}
\end{equation}
If exhaustion is established, the new bound is zero. Exhausted channels are skipped, and equal reductions are resolved by a fixed channel order (dense first). This rule also initializes reading when no document has been observed.

After each batch, update the bounds and extend the certified prefix. The batch size is capped by the remaining budget. When the budget is consumed, or both channels are exhausted, return the accumulated output. Certification depends on the observed bounds, so changing the read schedule affects the number of certified results without changing the meaning of an emitted position.
